\resizebox{\textwidth}{!}{%
\begin{tikzpicture}[every node/.style={font=\small}]

% Timeline axis
\draw[thick, -{Stealth[length=3mm]}] (-0.5,0) -- (16,0);

% Year markers
\foreach \x/\yr in {0.5/2020, 5/2021, 9.5/2022, 14/2024} {
  \draw[thick] (\x,0.15) -- (\x,-0.15);
  \node[below=0.3cm, font=\bfseries] at (\x,0) {\yr};
}

% Node boxes above timeline
\node[dia/lrbox=3.5cm, above=1cm of {(0.5,0)}] (n1) {%
  \textbf{Огляд інструментів WebAR}\\[2pt]
  A-Frame, AR.js, Three.js, JSARToolKit, Babylon.js\\[2pt]
%  {\footnotesize CEUR-WS Vol.~2832}%
};

\node[dia/lrbox=3.5cm, above=1cm of {(5,0)}] (n2) {%
  \textbf{Квести WebAR + методи дизайну ІОР}\\[2pt]
  маркерне та GPS-орієнтоване AR; методи проєктування\\[2pt]
%  {\footnotesize IOP JPhysConf; ACM DHW}%
};

\node[dia/lrbox=3.5cm, above=1cm of {(9.5,0)}] (n3) {%
  \textbf{6-компонентна модель проєктування ІОР}\\[2pt]
  аналіз~$\to$ вимоги~$\to$ ПЗ~$\to$ розробка~$\to$ впровадження~$\to$ оцінювання\\[2pt]
%  {\footnotesize Educational Dimension}%
};

\node[dia/lrbox=3.5cm, above=1cm of {(14,0)}] (n4) {%
  \textbf{Методика WebAR + ML}\\[2pt]
  TensorFlow.js, Teachable Machine, MindAR\\[2pt]
%  {\footnotesize CEUR-WS CoSinE; Educational Dimension}%
};

% Vertical connectors from timeline to boxes
\draw[dia/line] (0.5,0.15) -- (n1.south);
\draw[dia/line] (5,0.15) -- (n2.south);
\draw[dia/line] (9.5,0.15) -- (n3.south);
\draw[dia/line] (14,0.15) -- (n4.south);

% Horizontal progression arrows between boxes
\draw[dia/arrow] (n1.east) -- (n2.west);
\draw[dia/arrow] (n2.east) -- (n3.west);
\draw[dia/arrow] (n3.east) -- (n4.west);

% Bottom label: maturity progression
\node[below=1cm, font=\footnotesize\itshape] at (0.5,0) {Рівень 2};
\node[below=1cm, font=\footnotesize\itshape] at (5,0) {Рівень 2--3};
\node[below=1cm, font=\footnotesize\itshape] at (9.5,0) {Рівень 3};
\node[below=1cm, font=\footnotesize\itshape] at (14,0) {Рівень 3--4};

\end{tikzpicture}%
}
